\scanpage{28}\mspage{26}
la cat des foncteurs de $I$ dans $I'$, compatibles
avec source et but. Si cette catégorie est $\ne\varnothing$,
elle est $0$-connexe (et, je présume, même
contractile ?).

Je vais admettre provisoirement qu'il en est ainsi,
\inserted{du moins} dans le cas particulier de $\unsure{\mathscr{L}|\mathcal{L}}_4$ (intervalle gén. de long. $8$)
et $\mathord{\tikz[baseline=0pt,x=1ex,y=1ex]{\draw[line width=.09ex,line join=round] (0,0) -- (.95,1.30) -- (1.52,0) -- cycle; \draw[line width=.09ex] (.32,0) -- (1.02,1.12);}}_2$. Cela donne donc un foncteur
\[
\Theta\colon X\longrightarrow X_{\unsure{\Sigma|Z}}\quad .
\]
On a gagné, car la propriété universelle de ce
foncteur est pratiquement \unsure{évidente|immédiate}. \deleted{Cela établit} \inserted{Mais} j'ai
oublié de vérifier que $\Theta$ transforme \deleted{objets} \inserted{\unsure{élts|él.}} de $\Sigma$
\deleted{dans $X_\Sigma$}
en flèches inversibles. Soit donc
\diagram{\begin{tikzcd}[ampersand replacement=\&,column sep=1.15em]
x \arrow[r,"{u}"] \& y \arrow[r,<-] \& {}
\end{tikzcd}}
avec $u\in\Sigma$, on va définir un chemin en sens
inverse
\diagram{\begin{tikzcd}[ampersand replacement=\&,column sep=1.15em]
y \arrow[r,no head,"{c}"] \& x
\end{tikzcd}}
qui définira dans $X_\Sigma$ une flèche inverse de $u$.
L'idée naturelle est de prendre le ``chemin''
\diagram{\begin{tikzcd}[ampersand replacement=\&]\deleted{\unsure{x|y}}\&\deleted{\unsure{z|x}}\arrow[l,"{\deleted{\unsure{u}}}"]\end{tikzcd}}
\diagram{\begin{tikzcd}[ampersand replacement=\&,column sep=1.15em]
y \arrow[r,<-,"{u}"] \& x
\end{tikzcd}}
de $y$ dans $x$, mais ce n'est pas un chemin au sens
de nos définitions, on va le compléter en
\diagram{\begin{tikzcd}[ampersand replacement=\&,column sep=1.15em]
c\colon\quad y \arrow[r,"{\mathrm{id}}"] \& y \arrow[r,<-,"{u}"] \& x \arrow[r,no head,"{\mathrm{id}}"] \& x \arrow[r,<-,"{\mathrm{id}}"] \& x
\end{tikzcd}}
Je dis que $c$ fait l'affaire. on doit prouver
deux relations d'homotopie, de \deleted{$\unsure{c\ast u}$} $c.\bar{u}$ et
de $\bar{u}.c$ avec les chemins\deleted{\unsure{de}} identiques \unsure{en|au} $x$ de long.
$8$ (indice $4$ i.e. dans $\underline{\mathrm{Ch}}_{\Sigma,4}(X;x,\unsure{x|X})$ resp $\underline{\mathrm{Ch}}_{\Sigma,4}(X;y,y)$)
% archive sheet number: 27
